\documentclass[11pt]{article}
% Shared preamble for "The Kodamai Container Program: A Complete Technical Record"
% Self-contained; compiles with pdflatex.
\usepackage[utf8]{inputenc}
\usepackage[T1]{fontenc}
\usepackage{lmodern}
\usepackage{amsmath,amssymb,amsthm}
\usepackage{stmaryrd}   % \llbracket \rrbracket
\usepackage{mathrsfs}
\usepackage[margin=1in]{geometry}
\usepackage{enumitem}
\usepackage{microtype}
\usepackage{newunicodechar}
\usepackage[colorlinks=true,linkcolor=blue,citecolor=blue,urlcolor=blue]{hyperref}

% ---- Theorem environments -------------------------------------------------
\theoremstyle{plain}
\newtheorem{theorem}{Theorem}[section]
\newtheorem{proposition}[theorem]{Proposition}
\newtheorem{corollary}[theorem]{Corollary}
\newtheorem{lemma}[theorem]{Lemma}
\theoremstyle{definition}
\newtheorem{definition}[theorem]{Definition}
\newtheorem{result}[theorem]{Result}
\theoremstyle{remark}
\newtheorem*{remark}{Remark}
\newtheorem*{flag}{Flag}

% ---- Grade / registry helper macros ---------------------------------------
% Trust grades are copied verbatim from the master outline; never inflated.
\newcommand{\grade}[1]{\textbf{[#1]}}
\newcommand{\node}[1]{\texttt{#1}}

% ---- Unicode -> LaTeX mapping (so the outline prose transcribes verbatim) --
% Superscripts / subscripts
\newunicodechar{¹}{\ensuremath{^{1}}}
\newunicodechar{²}{\ensuremath{^{2}}}
\newunicodechar{³}{\ensuremath{^{3}}}
\newunicodechar{⁰}{\ensuremath{^{0}}}
\newunicodechar{⁺}{\ensuremath{^{+}}}
\newunicodechar{₁}{\ensuremath{_{1}}}
\newunicodechar{₂}{\ensuremath{_{2}}}
\newunicodechar{₃}{\ensuremath{_{3}}}
\newunicodechar{₄}{\ensuremath{_{4}}}
\newunicodechar{₊}{\ensuremath{_{+}}}
% Greek
\newunicodechar{Γ}{\ensuremath{\Gamma}}
\newunicodechar{Δ}{\ensuremath{\Delta}}
\newunicodechar{Π}{\ensuremath{\Pi}}
\newunicodechar{Σ}{\ensuremath{\Sigma}}
\newunicodechar{Ψ}{\ensuremath{\Psi}}
\newunicodechar{δ}{\ensuremath{\delta}}
\newunicodechar{ε}{\ensuremath{\varepsilon}}
\newunicodechar{θ}{\ensuremath{\theta}}
\newunicodechar{λ}{\ensuremath{\lambda}}
\newunicodechar{ν}{\ensuremath{\nu}}
\newunicodechar{φ}{\ensuremath{\varphi}}
\newunicodechar{ω}{\ensuremath{\omega}}
% Blackboard / script / fraktur
\newunicodechar{ℕ}{\ensuremath{\mathbb{N}}}
\newunicodechar{ℤ}{\ensuremath{\mathbb{Z}}}
\newunicodechar{ℰ}{\ensuremath{\mathcal{E}}}
\newunicodechar{𝒟}{\ensuremath{\mathcal{D}}}
\newunicodechar{𝒯}{\ensuremath{\mathcal{T}}}
\newunicodechar{𝔤}{\ensuremath{\mathfrak{g}}}
\newunicodechar{𝔽}{\ensuremath{\mathbb{F}}}
% Arrows and relations
\newunicodechar{→}{\ensuremath{\to}}
\newunicodechar{↔}{\ensuremath{\leftrightarrow}}
\newunicodechar{⟹}{\ensuremath{\implies}}
\newunicodechar{⟺}{\ensuremath{\iff}}
\newunicodechar{≃}{\ensuremath{\simeq}}
\newunicodechar{≅}{\ensuremath{\cong}}
\newunicodechar{≠}{\ensuremath{\neq}}
\newunicodechar{≡}{\ensuremath{\equiv}}
\newunicodechar{≤}{\ensuremath{\leq}}
\newunicodechar{≥}{\ensuremath{\geq}}
\newunicodechar{⊇}{\ensuremath{\supseteq}}
\newunicodechar{⊊}{\ensuremath{\subsetneq}}
\newunicodechar{∈}{\ensuremath{\in}}
% Operators
\newunicodechar{×}{\ensuremath{\times}}
\newunicodechar{∂}{\ensuremath{\partial}}
\newunicodechar{−}{\ensuremath{-}}
\newunicodechar{∘}{\ensuremath{\circ}}
\newunicodechar{∞}{\ensuremath{\infty}}
\newunicodechar{∧}{\ensuremath{\wedge}}
\newunicodechar{∨}{\ensuremath{\vee}}
\newunicodechar{⊕}{\ensuremath{\oplus}}
\newunicodechar{⊗}{\ensuremath{\otimes}}
\newunicodechar{⊙}{\ensuremath{\odot}}
\newunicodechar{⋆}{\ensuremath{\star}}
\newunicodechar{⋈}{\ensuremath{\bowtie}}
\newunicodechar{⋉}{\ensuremath{\ltimes}}
\newunicodechar{⋊}{\ensuremath{\rtimes}}
\newunicodechar{▷}{\ensuremath{\triangleright}}
\newunicodechar{◁}{\ensuremath{\triangleleft}}
\newunicodechar{⟦}{\ensuremath{\llbracket}}
\newunicodechar{⟧}{\ensuremath{\rrbracket}}
% Punctuation
\newunicodechar{–}{--}
\newunicodechar{—}{---}
\newunicodechar{§}{\S}


\title{The Kodamai Container Program: A Complete Technical Record\\[4pt]
\large Part V --- Formalisation in Lean 4}
\author{MacBeth}
\date{2026-10-10}

\begin{document}
\maketitle
\begin{center}\itshape Prepared for Neil Ghani.\end{center}

\begin{abstract}
Part~V records the Lean~4 formalisation. The core equivalence-chain links and the cohomological
obstruction engines are machine-checked: each oracle certifies a precise, honestly-scoped
statement; most are Mathlib-free (Lean~4 core) and axiom-minimal (\texttt{propext}/\texttt{Quot.sound}
only, several axiom-free via \texttt{decide}). Every item is \grade{lean-verified} unless stated
otherwise.
\end{abstract}

\paragraph{Thesis of Part V.}
The theorems are not hand-waving: the core equivalence-chain links and the cohomological obstruction
engines are machine-checked in Lean~4. Each oracle certifies a precise, honestly-scoped statement;
the general cohomology theory, where invoked, is cited as a black box (Brown VI.9), with the Lean
file certifying only a finite/explicit computation.

\section{The equivalence-chain milestones (M1--M4)}
\emph{Precis.} The container library and the directed-container $\leftrightarrow$ comonad
$\leftrightarrow$ small-category equivalence, plus the ZS associativity characterisation.

\begin{theorem}[M1/M2/M2b/M4 $+$ ZS assoc, \grade{lean-verified}]
\texttt{Basic}, \texttt{Directed}, \texttt{ComonadConverse}, \texttt{Cofunctor}, \texttt{DContCat},
\texttt{ZappaSzep} --- milestones M1/M2/M2b/M4 and ZS associativity (ZS1--ZS4 $\iff$ assoc).
\end{theorem}

\begin{theorem}[M3, \grade{lean-verified}]
\texttt{Comonoid.lean} / \texttt{ComonoidConverse.lean} --- directed container $\iff$
$\triangleleft$-comonoid (both directions).
\end{theorem}

\begin{theorem}[Cont and the four monoidal structures, \grade{lean-verified}]
\texttt{Cont.lean}, \texttt{Sequential.lean}, \texttt{Monoidal.lean}, \texttt{Dirichlet.lean},
\texttt{FourMonoidal.lean} --- Cont as a category, the $\triangleleft$ operator, all four monoidal
structures $+$ comparison isos.
\end{theorem}

\section{The free monad}
\begin{theorem}[Free monad, \grade{lean-verified}]
The free monad on a container $=$ a $\triangleleft$-monoid; the three monoid laws
(assoc $=$ \texttt{graft\_assoc}), \texttt{Quot.sound}-only. Node \node{free-monad-grafting}
(\texttt{Free.lean}); a partial universal property in \texttt{FreeUniversal.lean}. (The cofree
comonad is NOT core-formalisable --- it needs coinductive \texttt{PFunctor.M}.)
\end{theorem}

\section{The cohomological obstruction engines ($H^2$)}
\emph{Precis.} Three standalone, Mathlib-free Lean files compute the group cohomology that homes the
ZS and inventory obstructions. Honest scope: each certifies a FINITE/explicit computation; the
general cohomology theory is cited as a black box (Brown VI.9).

\begin{theorem}[Cyclic norm engine, \grade{lean-verified}]
$H^2(C_n;ℤ)\cong ℤ/n$ via the norm map. Node \node{cyclic-h2-norm-engine}.
\end{theorem}

\begin{theorem}[Fibre $ℤ/2$ engine, \grade{lean-verified}]
$H^2(ℤ/2;ℤ/2)\cong ℤ/2$, the fibre edge, with nonsplit-$ℤ/4$ vs.\ split-Klein witnesses. Node
\node{fibre-h2-z2-engine}.
\end{theorem}

\begin{theorem}[Unified gcd engine, \grade{lean-verified}]
$H^2(C_n;ℤ/m)\cong ℤ/\gcd(n,m)$, subsuming both engines above, with a from-scratch B\'ezout. Node
\node{cyclic-h2-gcd-engine}.
\end{theorem}

\section{The gerbe / reduced-nerve engine and $o_ν$ slices}
\emph{Precis.} Machine-checked witnesses for the supply-chain gerbe and the skew-brace solvability
obstruction.

\begin{theorem}[Reduced-nerve $𝔽_2$ gerbe engine, \grade{lean-verified}]
$H^2(B(ℤ/2))=1$, diamond $(1,1,0)$, $D_4$ non-split --- the three gerbe witnesses (MEMORY
\texttt{lean-gerbe-h2-witnesses}; grounds \node{supply-chain-inventory-gerbe-h2}).
\end{theorem}

\begin{theorem}[Nerve-level separation oracle, \grade{lean-verified}]
$N_1\subsetneq N_2$ machine-checked, bar complex $𝔽_2$, Betti $(1,1,1)$ vs.\ $(1,0,0)$ (grounds
\node{composition-nerve-level}).
\end{theorem}

\begin{theorem}[$o_ν$ finite slices, \grade{lean-verified}]
The $|G|=8$ additive slice, the $|G|=16$ three-way lock $+$ torsor dichotomy, and the corrected
$D=ℤ/8$ torsor criterion --- all \texttt{decide}-based, axiom-free or \texttt{propext}-only (nodes
under \node{solvability-obstruction-o-nu} / \node{quiver-skew-brace-zs}).
\end{theorem}

\begin{theorem}[Prunability descent cochain, \grade{lean-verified}]
Prunability $\iff δ\equiv0$ on the $|L|=2$ linearisation (node under \node{quiver-skew-brace-zs}).
\end{theorem}

\section{Transfer and distributivity cells}
\begin{theorem}[Transfer, \grade{lean-verified}]
Monad$\to$comonad transfer and its dual, three laws each (\texttt{MonadComonadTransfer.lean},
\texttt{DualTransfer.lean}).
\end{theorem}
\begin{theorem}[Closed Dirichlet and distributivity cells, \grade{lean-verified}]
Dirichlet closed / $Π$-form internal hom (axiom-free); the four Hedges distributivity cells
($\triangleleft/\times$, $\triangleleft/+$, $\otimes/+$, $\times/+$); and the Workers retract
$ΔS\otimes ΔT$ of $ΔS\triangleleft ΔT$.
\end{theorem}

\end{document}
